\section{Logging, observabilidad y auditoria}

\subsection{Descripcion}

Este documento define la relacion entre logs estructurados, observabilidad tecnica y auditoria funcional. Su objetivo es dejar claro que problemas resuelve cada uno, donde se complementan y como evitar duplicacion o lagunas de trazabilidad.

\subsection{Alcance}

Aplica a eventos tecnicos, correlacion operativa, rastreo de errores y acciones funcionales criticas del negocio. No define dashboards concretos ni una implementacion detallada de telemetria.

\subsection{Definiciones}

\begin{itemize}
  \item \textbf{Log estructurado}: evento tecnico emitido en formato consistente y parseable.
  \item \textbf{Observabilidad}: capacidad de entender el comportamiento del sistema a partir de señales tecnicas.
  \item \textbf{Auditoria}: registro funcional de acciones relevantes para negocio, operacion y trazabilidad.
  \item \textbf{Correlacion}: capacidad de relacionar multiples eventos tecnicos y funcionales dentro de una misma operacion.
\end{itemize}

\subsection{Reglas}

\begin{itemize}
  \item Logs, observabilidad y auditoria son complementarios y no deben confundirse.
  \item Un log tecnico no sustituye un evento de auditoria cuando existe impacto funcional relevante.
  \item Toda accion critica del sistema debe dejar al menos el rastro tecnico y funcional necesario para ser explicada.
  \item Los logs deben emitirse como flujo de eventos y no como archivos locales acoplados al proceso, en linea con Twelve-Factor App.
\end{itemize}

\subsubsection{Responsabilidad de cada capa}

\begin{itemize}
  \item Logs estructurados: diagnostico tecnico, errores, latencia, resultados de integracion y ejecucion de procesos.
  \item Observabilidad: metricas, trazas y correlacion para entender salud y comportamiento del sistema.
  \item Auditoria: evidencia funcional de acciones criticas como aprobacion de pagos, bloqueos, overrides, recalculos y archivado.
\end{itemize}

\subsubsection{Eventos funcionales que deben auditarse}

\begin{itemize}
  \item aprobacion o rechazo de comprobantes;
  \item bloqueo de usuarios;
  \item actualizacion manual de kickoff o estado de partido;
  \item carga o correccion manual de resultados;
  \item ejecucion de recalculo;
  \item archivado del torneo;
  \item cualquier accion administrativa que altere scoring, acceso o verdad oficial vigente.
\end{itemize}

\subsubsection{Campos minimos esperados}

\paragraph{Para logs tecnicos}

\begin{itemize}
  \item timestamp;
  \item nivel;
  \item modulo o servicio;
  \item tipo de evento;
  \item identificador de correlacion cuando exista;
  \item contexto tecnico suficiente para diagnostico sin filtrar secretos.
\end{itemize}

\paragraph{Para eventos de auditoria}

\begin{itemize}
  \item timestamp;
  \item actor;
  \item accion;
  \item entidad afectada;
  \item valor anterior y nuevo cuando aplique;
  \item justificacion o motivo cuando corresponda.
\end{itemize}

\subsubsection{Correlacion esperada}

\begin{itemize}
  \item Una operacion critica debe poder relacionar request, proceso interno, integracion externa y efecto funcional cuando aplique.
  \item Un error tecnico que cause impacto funcional debe ser rastreable hacia el proceso o evento funcional afectado.
  \item La auditoria debe permitir explicar por que un valor oficial cambio, mientras que la observabilidad debe ayudar a explicar como ocurrio tecnicamente.
\end{itemize}

\subsection{Edge Cases}

\begin{itemize}
  \item Una correccion de resultado puede requerir log tecnico, evento de auditoria y telemetria de recalculo sin que esos tres mecanismos repitan exactamente el mismo payload.
  \item Un error de integracion externa puede requerir alerta tecnica sin generar por si solo un evento de auditoria, salvo que derive en una accion manual relevante.
  \item Si una operacion administrativa falla antes de persistir cambios, puede requerir log tecnico y evidencia operativa sin convertirse en evento funcional exitoso.
\end{itemize}

\subsection{Invariantes}

\begin{itemize}
  \item Ningun evento critico queda sin rastro tecnico o funcional suficiente.
  \item La auditoria funcional no se usa como reemplazo de logs tecnicos.
  \item Los mecanismos de trazabilidad no deben exponer secretos ni datos innecesarios en texto plano.
\end{itemize}

\subsection{Preguntas abiertas}

\begin{itemize}
  \item Confirmar la estrategia final de correlacion entre logs, trazas y audit events, incluyendo si existira un identificador comun de operacion.
\end{itemize}

\subsection{Decisiones pendientes}

\begin{itemize}
  \item \texttt{Decision pendiente} \texttt{No bloqueante}: definir el catalogo final de eventos auditables y su esquema exacto una vez cierre la operacion administrativa real de v1.
  \item \texttt{Decision pendiente} \texttt{No bloqueante}: confirmar el nivel minimo de telemetria y alertas que se documentara formalmente en una iteracion posterior.
\end{itemize}
